\chapter{数据通路}

数据通路回答一个问题：值从哪里来，经过哪里，最后保存到哪里。前面已经有了寄存器和 ALU，但孤立部件还不能完成连续动作。必须把寄存器阵列、选择器、ALU、内部连线和写回路径接成一个可控制的闭环。

\section{一、寄存器阵列提供多个状态位置}

一个寄存器只能保存一个值。若硬件需要同时保存多个中间值，就需要寄存器阵列。寄存器阵列可以理解为一组编号寄存器，加上读写选择电路。

\begin{longtable}{p{0.22\linewidth}p{0.42\linewidth}p{0.24\linewidth}}
\toprule
端口 & 作用 & 类型 \\
\midrule
读地址 A/B & 选择两个源寄存器 & 输入 \\
读数据 A/B & 输出两个源值 & 输出 \\
写地址 & 选择目标寄存器 & 输入 \\
写数据 & 提供要保存的新值 & 输入 \\
写使能 & 决定是否在时钟沿写入 & 输入 \\
\bottomrule
\end{longtable}

读地址进入译码和选择网络，决定哪两个寄存器的值出现在读数据端口。写地址进入译码器，写使能有效时，只允许被选中的那个寄存器在时钟沿接收写数据。这样，一组寄存器就能被编号访问，而不是每个寄存器单独拉一套控制线。

\section{二、ALU 输入前必须选路}

ALU 的两个输入不一定总来自同一处。一个输入可能来自寄存器阵列，也可能来自计数器或固定 0；另一个输入可能来自寄存器、常量或某个扩展后的字段。因此 ALU 输入前通常有 mux。

\begin{lstlisting}
alu_a = mux(a_sel, reg_a, counter_value, zero)
alu_b = mux(b_sel, reg_b, constant, offset_value)
\end{lstlisting}

看一个具体动作：把寄存器 R1 和 R2 相加，结果写入 R3。数据通路需要让读地址 A 选择 R1，读地址 B 选择 R2；a\_sel 选择 reg\_a，b\_sel 选择 reg\_b；ALU op 选择 ADD。此时 ALU 输出就是 R1+R2。

若另一个动作是把某个计数器加 4，则 a\_sel 选择 counter\_value，b\_sel 选择 constant，ALU op 仍然可以是 ADD。同一个 ALU 因为输入选择不同，服务了不同硬件动作。

\section{三、写回路径决定结果保存到哪里}

ALU 输出不一定直接写回寄存器阵列。写回数据可能来自 ALU，也可能来自存储器读口、外部输入或计数器。写回前也需要 mux。

\begin{lstlisting}
write_data = mux(wb_sel, alu_result, memory_data, input_data)
\end{lstlisting}

继续看 R1+R2 写入 R3 的动作。wb\_sel 选择 alu\_result，写地址选择 R3，reg\_write 为 1。等时钟沿到来，R3 才真正保存结果。若 reg\_write 为 0，即使 ALU 算出了结果，寄存器阵列也不会改变。

这就是“写回”的硬件含义：某一路结果被选择出来，在明确时钟边界写入某个状态位置。结果不是因为算出来就自动保存，必须有写地址、写数据、写使能和时钟边界同时配合。

\section{四、内部连线不能被多个源同时强驱动}

当多个部件都可能把值送到同一个目的地时，不能让它们同时强行驱动同一条线。要么用 mux 明确选择一路，要么用总线协议规定同一时刻只有一个驱动者。

\begin{longtable}{p{0.22\linewidth}p{0.34\linewidth}p{0.32\linewidth}}
\toprule
方案 & 做法 & 风险 \\
\midrule
mux 选择 & 多路输入收成一路输出 & mux 层级太多会增加延迟 \\
共享总线 & 多个源共享数据线 & 必须避免多个源同时驱动 \\
\bottomrule
\end{longtable}

例如 ALU 输出、存储器输出和外部输入都想进入写回端口。如果三者直接接在一起，一旦两个源同时输出不同电平，就会形成冲突。mux 让这件事变成明确选择：当前周期只允许一路成为 write\_data。

\section{五、数据通路形成可计算闭环}

一个最小数据通路闭环可以写成：

\begin{lstlisting}
register file -> mux -> ALU -> mux -> register file
\end{lstlisting}

\begin{pdfdiagram}[x=1cm,y=1cm]
  \node[hw state,minimum width=2.0cm] (rf) at (1.3,1.55) {寄存器阵列};
  \node[hw compute,minimum width=1.3cm] (amux) at (3.8,2.1) {A mux};
  \node[hw compute,minimum width=1.3cm] (bmux) at (3.8,1.0) {B mux};
  \node[hw compute,minimum width=1.45cm] (alu) at (6.0,1.55) {ALU};
  \node[hw compute,minimum width=1.45cm] (wb) at (8.4,1.55) {写回 mux};
  \node[hw control,minimum width=1.8cm] (ctrl) at (5.1,3.05) {控制信号};
  \draw[hw bus] (rf.east) -- ++(0.65,0) |- (amux.west);
  \draw[hw bus] (rf.east) -- ++(0.65,0) |- (bmux.west);
  \draw[hw bus] (amux.east) -- (alu.west);
  \draw[hw bus] (bmux.east) -- (alu.west);
  \draw[hw bus] (alu.east) -- (wb.west);
  \draw[hw bus] (wb.east) -- ++(0.8,0) |- node[below,hw label] {时钟沿写入} (rf.south);
  \draw[hw bus] (ctrl.south) -- (amux.north);
  \draw[hw bus] (ctrl.south) -- (bmux.north);
  \draw[hw bus] (ctrl.south) -- (alu.north);
  \draw[hw bus] (ctrl.south) -- (wb.north);
  \node[hw label,align=center] at (5.1,0.15) {数据通路给出可走的路径，控制信号决定本周期哪条路径生效};
\end{pdfdiagram}
\diagramnote{最小数据通路是读旧状态、选输入、计算、选写回、在边界保存新状态的闭环。}

它能在一个周期内读取旧值、选择输入、完成组合计算，并在时钟沿保存结果。这个闭环已经能表达很多动作，但它还不会自己决定“当前周期做什么”。哪些地址送入寄存器阵列，ALU 做加法还是减法，写回选择哪一路，写使能是否打开，都需要另一个硬件部件统一产生。下一章就是控制器。

\begin{classicnote}
数据通路的标注方式是 trace：从哪个状态元件读出旧值，经过哪些 mux 和组合部件，到哪个状态元件写入新值。

\begin{itemize}
\item 纸上实验：为 \code{R3 = R1 + R2} 写一行 trace：读地址、ALU 输入选择、\code{alu\_op}、写回选择、写地址、写使能。
\item 机器证据：寄存器读口通常是组合输出，写口通常在时钟沿提交；读旧值和写新值不能混成同一个瞬间。
\item 冲突证据：内部总线或共享线同一时刻只能由一个源驱动；多源并联必须通过 mux、三态规则或仲裁结构解决。
\item 检查题：若写回 mux 选择线接错，ALU 算对了但寄存器结果错。排查时应看 ALU 输出，还是写回路径和写使能？
\end{itemize}

经典教材常把每条机器指令拆成数据通路上的值流。这里还没有完整指令集，但同样的方法已经可用：看值从哪里来，到哪里去，哪条控制线让它走这条路。
\end{classicnote}
